% --- Archon physics formalization source begin ---
% archon:physics
% archon:covers PhyXMiniProblems/problem_phyx_mini_0343.lean
% archon:source-report reports/phyx_mini/problem_phyx_mini_0343.source.json

\chapter{Physics problem phyx\_mini\_0343}
\label{ch:PhyXMiniProblems_problem_phyx_mini_0343}

\paragraph{Problem source.}
\#\# Physical scenario

\textbackslash{}textbf\{Figure\} shows a \$pV\$-diagram for \$0.0040\textasciitilde{}\textbackslash{}text\{mol\}\$ of \textbackslash{}textit\{ideal\} H\$\_2\$ gas. The temperature of the gas does not change during segment \$bc\$.

\#\# Question

What volume does this gas occupy at point \$c\$?

\#\# Answer choices

- A: A: 0.83L

- B: B: 0.8L

- C: C: 0.9L

- D: D: 1.8L

\#\# Figure caption (auxiliary; use the image as primary evidence)

The image is a pressure-volume (p-V) graph with the following features:

- Axes: The vertical axis is labeled "p (atm)" indicating pressure in atmospheres, and the horizontal axis is labeled "V (L)" indicating volume in liters.
- Points: There are three points labeled \textbackslash{}(a\textbackslash{}), \textbackslash{}(b\textbackslash{}), and \textbackslash{}(c\textbackslash{}).
  - Point \textbackslash{}(a\textbackslash{}) is located at (0.20 L, 0.50 atm).
  - Point \textbackslash{}(b\textbackslash{}) is directly above point \textbackslash{}(a\textbackslash{}) at (0.20 L, 2.0 atm).
  - Point \textbackslash{}(c\textbackslash{}) is to the right of point \textbackslash{}(a\textbackslash{}) at a lower pressure (0.50 atm) but higher volume.
- Lines and Arrows: 
  - A vertical line connects point \textbackslash{}(a\textbackslash{}) to point \textbackslash{}(b\textbackslash{}).
  - A curved line connects point \textbackslash{}(b\textbackslash{}) to point \textbackslash{}(c\textbackslash{}).
  - A horizontal line connects point \textbackslash{}(c\textbackslash{}) back to point \textbackslash{}(a\textbackslash{}).
- Arrows indicate the direction of a cyclic process:
  - From \textbackslash{}(a\textbackslash{}) to \textbackslash{}(b\textbackslash{}) is upward.
  - From \textbackslash{}(b\textbackslash{}) to \textbackslash{}(c\textbackslash{}) is downward along the curve.
  - From \textbackslash{}(c\textbackslash{}) to \textbackslash{}(a\textbackslash{}) is leftward along the horizontal line.

The graph shows a closed cycle process in a thermodynamic system.

\#\# Dataset metadata

Ideal Gases and Kinetic Theory; Physical Model Grounding Reasoning, Spatial Relation Reasoning

\paragraph{Recorded answer/context.}
B: B: 0.8L

\paragraph{Figure/image path.}
/root/proposal\_for\_physic/science-mango/phyx\_mini\_run/phyx\_data/test\_image/343.png

\paragraph{Formalization target.}
create a compiling Lean file with sorry bodies at `PhyXMiniProblems/problem\_phyx\_mini\_0343.lean`. The Lean declarations must preserve the physical quantities, dimensions or dimensional roles, figure labels, governing-law hypotheses, and final relation expressed by this problem.
Use Mathlib/Physlib names found through LeanExplore where available. If a physics API is missing, introduce faithful local abstractions rather than scalar placeholder aliases.

\begin{theorem}[Physics formalization target]
\label{thm:physics:phyx_mini_0343:target}
\lean{PhyXMiniProblems.ProblemPhyXMini0343.volumeAtPointC_matches_recordedAnswerB}
\uses{def:physics:phyx-mini-0343:phyxminiproblems-problemphyxmini0343-volumeinliters, def:physics:phyx-mini-0343:phyxminiproblems-problemphyxmini0343-idealhydrogenpvcycle, def:physics:phyx-mini-0343:phyxminiproblems-problemphyxmini0343-matchesproblemandfigure, def:physics:phyx-mini-0343:phyxminiproblems-problemphyxmini0343-hasphysicalidealgasparameters, def:physics:phyx-mini-0343:phyxminiproblems-problemphyxmini0343-satisfiesmolaridealgaslaw, def:physics:phyx-mini-0343:phyxminiproblems-problemphyxmini0343-recordedanswerchoice}
The assigned autoformalize agent should translate this physics problem into Lean declarations in the covered file, with theorem and lemma proof bodies written as `by sorry`.
\end{theorem}
\begin{proof}
This is an autoformalization task, not a proof task. Produce statements that can later be proved by the physics prover without weakening the physical model.
\end{proof}

% --- Archon declaration topology begin ---
\section{Declaration topology}
\begin{definition}[Pressure Quantity]
  \label{def:physics:phyx-mini-0343:phyxminiproblems-problemphyxmini0343-pressurequantity}
  \lean{PhyXMiniProblems.ProblemPhyXMini0343.PressureQuantity}
  Physical pressure, represented by Physlib's dimensionful pressure type
\end{definition}
\begin{proof}
  This is the declared model definition of Pressure Quantity.
\end{proof}
\begin{definition}[Volume Quantity]
  \label{def:physics:phyx-mini-0343:phyxminiproblems-problemphyxmini0343-volumequantity}
  \lean{PhyXMiniProblems.ProblemPhyXMini0343.VolumeQuantity}
  A nonnegative physical volume with dimension length \textasciicircum{} 3
\end{definition}
\begin{proof}
  This is the declared model definition of Volume Quantity.
\end{proof}
\begin{definition}[Temperature Quantity]
  \label{def:physics:phyx-mini-0343:phyxminiproblems-problemphyxmini0343-temperaturequantity}
  \lean{PhyXMiniProblems.ProblemPhyXMini0343.TemperatureQuantity}
  Absolute thermodynamic temperature
\end{definition}
\begin{proof}
  This is the declared model definition of Temperature Quantity.
\end{proof}
\begin{definition}[pressure In Pascals]
  \label{def:physics:phyx-mini-0343:phyxminiproblems-problemphyxmini0343-pressureinpascals}
  \lean{PhyXMiniProblems.ProblemPhyXMini0343.pressureInPascals}
  \uses{def:physics:phyx-mini-0343:phyxminiproblems-problemphyxmini0343-pressurequantity}
  Coherent-SI pressure readout in pascals
\end{definition}
\begin{proof}
  This is the declared model definition of pressure In Pascals.
\end{proof}
\begin{definition}[pressure In Atmospheres]
  \label{def:physics:phyx-mini-0343:phyxminiproblems-problemphyxmini0343-pressureinatmospheres}
  \lean{PhyXMiniProblems.ProblemPhyXMini0343.pressureInAtmospheres}
  \uses{def:physics:phyx-mini-0343:phyxminiproblems-problemphyxmini0343-pressurequantity, def:physics:phyx-mini-0343:phyxminiproblems-problemphyxmini0343-pressureinpascals}
  Pressure readout in standard atmospheres. The denominator is Physlib's dimensionful standard-atmosphere constant, rather than an untyped conversion factor hidden in the model
\end{definition}
\begin{proof}
  This is the declared model definition of pressure In Atmospheres.
\end{proof}
\begin{definition}[volume In Liters]
  \label{def:physics:phyx-mini-0343:phyxminiproblems-problemphyxmini0343-volumeinliters}
  \lean{PhyXMiniProblems.ProblemPhyXMini0343.volumeInLiters}
  \uses{def:physics:phyx-mini-0343:phyxminiproblems-problemphyxmini0343-volumequantity}
  Volume readout in litres; one cubic metre is one thousand litres
\end{definition}
\begin{proof}
  This is the declared model definition of volume In Liters.
\end{proof}
\begin{definition}[temperature Readout]
  \label{def:physics:phyx-mini-0343:phyxminiproblems-problemphyxmini0343-temperaturereadout}
  \lean{PhyXMiniProblems.ProblemPhyXMini0343.temperatureReadout}
  \uses{def:physics:phyx-mini-0343:phyxminiproblems-problemphyxmini0343-temperaturequantity}
  Numerical readout of absolute temperature in the common temperature unit used by the molar gas constant below. Only equality of the b and c readings is needed by this problem
\end{definition}
\begin{proof}
  This is the declared model definition of temperature Readout.
\end{proof}
\begin{definition}[Chemical Species]
  \label{def:physics:phyx-mini-0343:phyxminiproblems-problemphyxmini0343-chemicalspecies}
  \lean{PhyXMiniProblems.ProblemPhyXMini0343.ChemicalSpecies}
  Chemical species named in the problem statement
\end{definition}
\begin{proof}
  This is the declared model definition of Chemical Species.
\end{proof}
\begin{definition}[Equation Of State Model]
  \label{def:physics:phyx-mini-0343:phyxminiproblems-problemphyxmini0343-equationofstatemodel}
  \lean{PhyXMiniProblems.ProblemPhyXMini0343.EquationOfStateModel}
  Equation-of-state model specified by the word “ideal”
\end{definition}
\begin{proof}
  This is the declared model definition of Equation Of State Model.
\end{proof}
\begin{definition}[Gas Sample]
  \label{def:physics:phyx-mini-0343:phyxminiproblems-problemphyxmini0343-gassample}
  \lean{PhyXMiniProblems.ProblemPhyXMini0343.GasSample}
  \uses{def:physics:phyx-mini-0343:phyxminiproblems-problemphyxmini0343-chemicalspecies, def:physics:phyx-mini-0343:phyxminiproblems-problemphyxmini0343-equationofstatemodel}
  The physical gas sample. amountInMoles is an explicitly unit-labelled scalar readout, not a replacement type for the sample itself
\end{definition}
\begin{proof}
  This is the declared model definition of Gas Sample.
\end{proof}
\begin{definition}[Diagram Point]
  \label{def:physics:phyx-mini-0343:phyxminiproblems-problemphyxmini0343-diagrampoint}
  \lean{PhyXMiniProblems.ProblemPhyXMini0343.DiagramPoint}
  The three labelled thermodynamic states in the primary figure
\end{definition}
\begin{proof}
  This is the declared model definition of Diagram Point.
\end{proof}
\begin{definition}[Diagram Segment]
  \label{def:physics:phyx-mini-0343:phyxminiproblems-problemphyxmini0343-diagramsegment}
  \lean{PhyXMiniProblems.ProblemPhyXMini0343.DiagramSegment}
  The three directed process segments in cycle order
\end{definition}
\begin{proof}
  This is the declared model definition of Diagram Segment.
\end{proof}
\begin{definition}[start Point]
  \label{def:physics:phyx-mini-0343:phyxminiproblems-problemphyxmini0343-diagramsegment-startpoint}
  \lean{PhyXMiniProblems.ProblemPhyXMini0343.DiagramSegment.startPoint}
  \uses{def:physics:phyx-mini-0343:phyxminiproblems-problemphyxmini0343-diagrampoint, def:physics:phyx-mini-0343:phyxminiproblems-problemphyxmini0343-diagramsegment}
  Initial endpoint of each arrow in the depicted cycle
\end{definition}
\begin{proof}
  This is the declared model definition of start Point.
\end{proof}
\begin{definition}[end Point]
  \label{def:physics:phyx-mini-0343:phyxminiproblems-problemphyxmini0343-diagramsegment-endpoint}
  \lean{PhyXMiniProblems.ProblemPhyXMini0343.DiagramSegment.endPoint}
  \uses{def:physics:phyx-mini-0343:phyxminiproblems-problemphyxmini0343-diagrampoint, def:physics:phyx-mini-0343:phyxminiproblems-problemphyxmini0343-diagramsegment}
  Final endpoint of each arrow in the depicted cycle
\end{definition}
\begin{proof}
  This is the declared model definition of end Point.
\end{proof}
\begin{definition}[Segment Shape]
  \label{def:physics:phyx-mini-0343:phyxminiproblems-problemphyxmini0343-segmentshape}
  \lean{PhyXMiniProblems.ProblemPhyXMini0343.SegmentShape}
  Qualitative shapes of process segments shown in the bitmap
\end{definition}
\begin{proof}
  This is the declared model definition of Segment Shape.
\end{proof}
\begin{definition}[Thermodynamic State]
  \label{def:physics:phyx-mini-0343:phyxminiproblems-problemphyxmini0343-thermodynamicstate}
  \lean{PhyXMiniProblems.ProblemPhyXMini0343.ThermodynamicState}
  \uses{def:physics:phyx-mini-0343:phyxminiproblems-problemphyxmini0343-pressurequantity, def:physics:phyx-mini-0343:phyxminiproblems-problemphyxmini0343-volumequantity, def:physics:phyx-mini-0343:phyxminiproblems-problemphyxmini0343-temperaturequantity}
  Pressure, volume, and temperature at one labelled state
\end{definition}
\begin{proof}
  This is the declared model definition of Thermodynamic State.
\end{proof}
\begin{definition}[Ideal Hydrogen PVCycle]
  \label{def:physics:phyx-mini-0343:phyxminiproblems-problemphyxmini0343-idealhydrogenpvcycle}
  \lean{PhyXMiniProblems.ProblemPhyXMini0343.IdealHydrogenPVCycle}
  \uses{def:physics:phyx-mini-0343:phyxminiproblems-problemphyxmini0343-gassample, def:physics:phyx-mini-0343:phyxminiproblems-problemphyxmini0343-diagrampoint, def:physics:phyx-mini-0343:phyxminiproblems-problemphyxmini0343-diagramsegment, def:physics:phyx-mini-0343:phyxminiproblems-problemphyxmini0343-segmentshape, def:physics:phyx-mini-0343:phyxminiproblems-problemphyxmini0343-thermodynamicstate}
  The gas, its three states, and the qualitative geometry of the displayed cycle. The gas constant is a readout in litre-atmospheres per mole per the temperature unit used by temperatureReadout
\end{definition}
\begin{proof}
  This is the declared model definition of Ideal Hydrogen PVCycle.
\end{proof}
\begin{definition}[Matches Problem And Figure]
  \label{def:physics:phyx-mini-0343:phyxminiproblems-problemphyxmini0343-matchesproblemandfigure}
  \lean{PhyXMiniProblems.ProblemPhyXMini0343.MatchesProblemAndFigure}
  \uses{def:physics:phyx-mini-0343:phyxminiproblems-problemphyxmini0343-pressureinatmospheres, def:physics:phyx-mini-0343:phyxminiproblems-problemphyxmini0343-volumeinliters, def:physics:phyx-mini-0343:phyxminiproblems-problemphyxmini0343-idealhydrogenpvcycle}
  Data stated in the prose or read directly from the primary bitmap. Point c's pressure and its qualitative position are included, but its requested volume is deliberately not assigned a numerical value here
\end{definition}
\begin{proof}
  This is the declared model definition of Matches Problem And Figure.
\end{proof}
\begin{definition}[Has Physical Ideal Gas Parameters]
  \label{def:physics:phyx-mini-0343:phyxminiproblems-problemphyxmini0343-hasphysicalidealgasparameters}
  \lean{PhyXMiniProblems.ProblemPhyXMini0343.HasPhysicalIdealGasParameters}
  \uses{def:physics:phyx-mini-0343:phyxminiproblems-problemphyxmini0343-pressureinatmospheres, def:physics:phyx-mini-0343:phyxminiproblems-problemphyxmini0343-volumeinliters, def:physics:phyx-mini-0343:phyxminiproblems-problemphyxmini0343-temperaturereadout, def:physics:phyx-mini-0343:phyxminiproblems-problemphyxmini0343-diagrampoint, def:physics:phyx-mini-0343:phyxminiproblems-problemphyxmini0343-idealhydrogenpvcycle}
  Positivity conditions selecting physically meaningful gas states
\end{definition}
\begin{proof}
  This is the declared model definition of Has Physical Ideal Gas Parameters.
\end{proof}
\begin{definition}[Satisfies Molar Ideal Gas Law]
  \label{def:physics:phyx-mini-0343:phyxminiproblems-problemphyxmini0343-satisfiesmolaridealgaslaw}
  \lean{PhyXMiniProblems.ProblemPhyXMini0343.SatisfiesMolarIdealGasLaw}
  \uses{def:physics:phyx-mini-0343:phyxminiproblems-problemphyxmini0343-pressureinatmospheres, def:physics:phyx-mini-0343:phyxminiproblems-problemphyxmini0343-volumeinliters, def:physics:phyx-mini-0343:phyxminiproblems-problemphyxmini0343-temperaturereadout, def:physics:phyx-mini-0343:phyxminiproblems-problemphyxmini0343-diagrampoint, def:physics:phyx-mini-0343:phyxminiproblems-problemphyxmini0343-idealhydrogenpvcycle}
  The molar ideal-gas equation p V = n R T, stated at every labelled state in the units used by the figure. This is a general governing law and contains no numerical value for the requested volume at c
\end{definition}
\begin{proof}
  This is the declared model definition of Satisfies Molar Ideal Gas Law.
\end{proof}
\begin{lemma}[pressure Volume Product eq along BC]
  \label{lem:physics:phyx-mini-0343:phyxminiproblems-problemphyxmini0343-pressurevolumeproduct-eq-alongbc}
  \lean{PhyXMiniProblems.ProblemPhyXMini0343.pressureVolumeProduct_eq_alongBC}
  \uses{def:physics:phyx-mini-0343:phyxminiproblems-problemphyxmini0343-pressureinatmospheres, def:physics:phyx-mini-0343:phyxminiproblems-problemphyxmini0343-volumeinliters, def:physics:phyx-mini-0343:phyxminiproblems-problemphyxmini0343-idealhydrogenpvcycle, def:physics:phyx-mini-0343:phyxminiproblems-problemphyxmini0343-matchesproblemandfigure, def:physics:phyx-mini-0343:phyxminiproblems-problemphyxmini0343-hasphysicalidealgasparameters, def:physics:phyx-mini-0343:phyxminiproblems-problemphyxmini0343-satisfiesmolaridealgaslaw}
  Along the isothermal segment b c, the common sample obeying the ideal-gas law has equal pressure--volume products at the endpoints
\end{lemma}
\begin{proof}
  The conclusion follows from the governing relations and figure readouts named in the statement of pressure Volume Product eq along BC.
\end{proof}
\begin{definition}[Answer Choice]
  \label{def:physics:phyx-mini-0343:phyxminiproblems-problemphyxmini0343-answerchoice}
  \lean{PhyXMiniProblems.ProblemPhyXMini0343.AnswerChoice}
  Labels of the four displayed answer choices
\end{definition}
\begin{proof}
  This is the declared model definition of Answer Choice.
\end{proof}
\begin{definition}[liters]
  \label{def:physics:phyx-mini-0343:phyxminiproblems-problemphyxmini0343-answerchoice-liters}
  \lean{PhyXMiniProblems.ProblemPhyXMini0343.AnswerChoice.liters}
  \uses{def:physics:phyx-mini-0343:phyxminiproblems-problemphyxmini0343-answerchoice}
  Volume in litres printed beside each answer label
\end{definition}
\begin{proof}
  This is the declared model definition of liters.
\end{proof}
\begin{definition}[recorded Answer Choice]
  \label{def:physics:phyx-mini-0343:phyxminiproblems-problemphyxmini0343-recordedanswerchoice}
  \lean{PhyXMiniProblems.ProblemPhyXMini0343.recordedAnswerChoice}
  \uses{def:physics:phyx-mini-0343:phyxminiproblems-problemphyxmini0343-answerchoice}
  The answer label recorded in the supplied dataset metadata
\end{definition}
\begin{proof}
  This is the declared model definition of recorded Answer Choice.
\end{proof}
% --- Archon declaration topology end ---
% --- Archon physics formalization source end ---
